\documentclass[11pt]{article}
\usepackage{amsmath,amsthm,amssymb}
\usepackage[margin=1in]{geometry}

\newtheorem{theorem}{Theorem}
\newtheorem{lemma}{Lemma}
\newtheorem{proposition}{Proposition}
\newtheorem{corollary}{Corollary}
\theoremstyle{definition}
\newtheorem{definition}{Definition}
\theoremstyle{remark}
\newtheorem{remark}{Remark}

\DeclareMathOperator{\Des}{Des}
\DeclareMathOperator{\SYT}{SYT}
\DeclareMathOperator{\ord}{ord}
\DeclareMathOperator{\supp}{supp}
\DeclareMathOperator{\core}{core}
\newcommand{\ZZ}{\mathbb{Z}}
\newcommand{\QQ}{\mathbb{Q}}
\newcommand{\la}{\lambda}
\newcommand{\inner}[1]{\left\langle #1 \right\rangle}
\newcommand{\zt}{\zeta_3}

\title{The graded order law collapses at $\zeta_3$:\\
$G_\la(\zeta_3)=0$ if and only if $\la\in\{(3,1),(2,1,1)\}$}
\author{Clio}
\date{2026-06-02}

\begin{document}
\maketitle

\begin{abstract}
Let $G_\la(q)=\sum_{T\in\SYT(\la)}q^{s(T)}$ be the graded branch--exponent generating
function, where $s(T)=\sum_{i\in\Des(T)}w_i$ with the parity weights $w_i=2i-1$ if
$n-i$ is odd and $w_i=0$ otherwise. At $q=-1$ this polynomial obeys a rich vanishing
law, $\ord_{q=-1}G_\la=\lfloor|2\text{-}\core(\la)|/2\rfloor$, unbounded along the
staircase of $2$-cores. We prove that at a primitive \emph{cube} root of unity $\zt$
the law collapses completely:
\[
   G_\la(\zt)=0 \iff \la\in\{(3,1),(2,1,1)\},
\]
and at those two shapes the order is exactly $1$. In particular every partition of
positive $3$-weight (every non-$3$-core, and indeed all but two $3$-cores) satisfies
$\ord_{\zt}G_\la=0$. The proof turns the apparent root-of-unity cancellation into a
question of Schur positivity. A master identity expresses the fiber value as
$G_\la(\zt)=\inner{s_\la,\,h_1^{e}\prod_{k}(h_2+e_2\zt^{2k-1})}$; grouping the factors
by residue class mod $3$ and using the single quadratic identity
$\Theta:=(h_2+e_2\zt)(h_2+e_2\zt^{2})=s_{(4)}+2s_{(2,2)}+s_{(1^4)}$ (manifestly
Schur-positive) reduces the value, in all cases, to a nonnegative integer inner
product $\inner{s_\la,h_1^{A}\Theta^{b}}$. A Pieri computation and the fact that
$\Theta^{b}$ has \emph{full} Schur support for $b\ge3$ finish the proof, the only
residue being a finite check up to $n=17$.
\end{abstract}

\section{Setup}

Throughout $\la\vdash n$, $\SYT(\la)$ is the set of standard Young tableaux of shape
$\la$, and for $T\in\SYT(\la)$,
\[
  \Des(T)=\{\,i\in[n-1]: i{+}1 \text{ lies in a strictly lower row of }T\text{ than }i\,\}.
\]
Fix the \emph{active set} and the parity weights
\[
  A=\{\,k\in[n-1]: n-k\text{ odd}\,\},\qquad
  w_k=(2k-1)\,[\,k\in A\,],\qquad
  s(T)=\sum_{i\in\Des(T)}w_i,
\]
and set $m=|A|=\lfloor n/2\rfloor$, $e=n\bmod 2$. The object of study is
\[
  G_\la(q)=\sum_{T\in\SYT(\la)}q^{s(T)}\ \in\ \ZZ_{\ge0}[q].
\]
We write $\zt=e^{2\pi i/3}$ for a primitive cube root of unity, $\Phi_3(q)=q^2+q+1$
for the third cyclotomic polynomial, and
$\ord_{\zt}G_\la=\max\{m:\Phi_3^m\mid G_\la\}$ for the order of vanishing at $\zt$.

We work in the ring $\Lambda$ of symmetric functions, with Schur, complete, elementary
and power-sum bases $s_\la,h_r,e_r,p_r$, and the Hall inner product
$\inner{\,\cdot\,,\cdot}$, for which $\inner{s_\la,s_\mu}=\delta_{\la\mu}$,
$\inner{s_\la,h_\mu}=K_{\la\mu}$ (Kostka numbers), and $\inner{s_\la,p_\mu}=\chi^\la(\mu)$.
For symmetric functions $f$ we write $\supp(f)=\{\mu:\inner{s_\mu,f}\neq0\}$; if all
Schur coefficients of $f$ are $\ge0$ we call $f$ \emph{Schur-positive}. We use freely
$h_2=\tfrac12(p_1^2+p_2)$, $e_2=\tfrac12(p_1^2-p_2)$, so that $h_2+e_2=p_1^2=h_1^2$ and
$h_2-e_2=p_2$.

\section{The master identity for the fiber value}

The following is the value-at-$\zt$ specialization of the descent enumerator that drove
the $q=-1$ law; the proof is identical in form, so we give it compactly.

\begin{lemma}[Descent enumerator; Clio 2026-06-01]\label{lem:descent}
For indeterminates $\eta=(\eta_1,\dots,\eta_{n-1})$, with $\alpha(R)$ the composition of
$n$ whose descent set is $R$,
\[
  \sum_{T\in\SYT(\la)}\ \prod_{k\in\Des(T)}\eta_k
   \;=\;\sum_{R\subseteq[n-1]} K_{\la,\alpha(R)}\ \prod_{k\in R}\eta_k\prod_{k\notin R}(1-\eta_k).
\]
\end{lemma}

\begin{proposition}[Master identity]\label{prop:master}
For any complex number $\zeta$,
\[
  G_\la(\zeta)\;=\;\Big\langle s_\la,\; h_1^{\,e}\prod_{k\in A}\big(h_2+e_2\,\zeta^{\,2k-1}\big)\Big\rangle .
\]
\end{proposition}

\begin{proof}
Put $\eta_k=\zeta^{\,w_k}$ in Lemma~\ref{lem:descent}. Then
$\prod_{k\in\Des(T)}\eta_k=\zeta^{\sum_{i\in\Des(T)}w_i}=\zeta^{s(T)}$, so the left side
is $G_\la(\zeta)$. For $k\notin A$ we have $w_k=0$, hence $\eta_k=1$ and the factor
$1-\eta_k$ vanishes; only $R\supseteq A^c$ survive. Write $R=A^c\sqcup U$ with
$U\subseteq A$; for $k\in A$, $\eta_k=\zeta^{2k-1}$ and $1-\eta_k=1-\zeta^{2k-1}$. As in
the $q=-1$ proof, the content of $\alpha(A^c\cup U)$ is $(2^{\,m-|U|},1^{\,e+2|U|})$
(adjoining $k\in A$ to the marks of $\alpha(A^c)$ splits the $2$-gap centred at $k$ into
two $1$'s), whence $K_{\la,\alpha(A^c\cup U)}=\inner{s_\la,h_2^{\,m-|U|}h_1^{\,e+2|U|}}$.
Therefore
\[
  G_\la(\zeta)=\Big\langle s_\la,\ h_1^{e}\sum_{U\subseteq A}
    h_2^{\,m-|U|}h_1^{\,2|U|}\prod_{k\in U}\zeta^{2k-1}\prod_{k\in A\setminus U}(1-\zeta^{2k-1})\Big\rangle .
\]
The sum over $U$ factorises: each $k\in A$ contributes $h_1^2\,\zeta^{2k-1}$ if $k\in U$
and $h_2\,(1-\zeta^{2k-1})$ if $k\notin U$, so the bracket is
$h_1^{e}\prod_{k\in A}\big(h_2(1-\zeta^{2k-1})+h_1^2\zeta^{2k-1}\big)$. Finally
$h_2(1-\zeta^{2k-1})+h_1^2\zeta^{2k-1}=h_2+(h_1^2-h_2)\zeta^{2k-1}=h_2+e_2\zeta^{2k-1}$.
\end{proof}

At $\zeta=-1$ each factor is $h_2-e_2=p_2$ and one recovers $G_\la(-1)=\chi^\la(2^m1^e)$;
the cube-root case is governed instead by the following grouping.

\section{Grouping by residue and the quadratic identity}

For $\zeta=\zt$ write $z_k=\zt^{\,2k-1}$. Since $2k-1\equiv 0,1,2\pmod 3$ according as
$k\equiv 2,1,0\pmod 3$, the value $z_k$ is $1,\zt,\zt^2$ respectively. Set
\[
  n_0=\#\{k\in A:k\equiv2\},\quad
  n_1=\#\{k\in A:k\equiv1\},\quad
  n_2=\#\{k\in A:k\equiv0\}\pmod 3,
\]
so that $n_0+n_1+n_2=m$, and define $\phi=h_2+e_2\zt$, $\bar\phi=h_2+e_2\zt^{2}$. Using
$h_2+e_2=h_1^2$,
\begin{equation}\label{eq:grouped}
  \prod_{k\in A}(h_2+e_2 z_k)=(h_1^2)^{\,n_0}\,\phi^{\,n_1}\,\bar\phi^{\,n_2}.
\end{equation}

\begin{lemma}[The imbalance is at most one]\label{lem:delta}
$\delta:=n_1-n_2\in\{0,1\}$, with $\delta=1$ if and only if $n\equiv 2\pmod 3$.
\end{lemma}

\begin{proof}
$A$ is an arithmetic progression of common difference $2$: it is
$\{1,3,\dots,n-1\}$ if $n$ is even and $\{2,4,\dots,n-1\}$ if $n$ is odd. Replacing $n$
by $n+6$ preserves the parity of $n$, hence the progression type, and appends to $A$
exactly the three new terms $x,x+2,x+4$ of the correct parity lying in $[n,n+5]$. These
three are consecutive terms of a step-$2$ progression, so modulo $3$ (as $\gcd(2,3)=1$)
they realise all three residues once each; thus each of $n_0,n_1,n_2$ increases by $1$
and $\delta$ is unchanged. Hence $\delta$ depends only on $n\bmod 6$, and a direct check
of $n=1,\dots,6$ gives $\delta=(0,1,0,0,1,0)$, i.e. $\delta=1\iff n\equiv2,5\pmod 6\iff
n\equiv2\pmod 3$.
\end{proof}

\begin{lemma}[Quadratic Schur-positivity]\label{lem:theta}
$\displaystyle \Theta:=\phi\,\bar\phi=h_2^2-h_2e_2+e_2^2=\frac{p_1^4+3p_2^2}{4}
   =s_{(4)}+2\,s_{(2,2)}+s_{(1^4)}.$
\end{lemma}

\begin{proof}
$\phi\bar\phi=h_2^2+h_2e_2(\zt+\zt^2)+e_2^2=h_2^2-h_2e_2+e_2^2$ since $\zt+\zt^2=-1$.
Substituting $h_2=\tfrac12(p_1^2+p_2)$, $e_2=\tfrac12(p_1^2-p_2)$ gives
$\tfrac14(p_1^4+3p_2^2)$. For the Schur expansion use Jacobi--Trudi /
Pieri: $h_2^2=s_{4}+s_{31}+s_{22}$, $e_2^2=s_{22}+s_{211}+s_{1^4}$,
$h_2e_2=s_{31}+s_{211}$; hence $h_2^2-h_2e_2+e_2^2=s_{4}+2s_{22}+s_{1^4}$.
\end{proof}

\section{Reduction to a Schur-support statement}

Combining Proposition~\ref{prop:master}, \eqref{eq:grouped} and
Lemmas~\ref{lem:delta}--\ref{lem:theta}, set
\[
  a:=e+2n_0\ \ (\ge0),\qquad b:=n_2\ \ (\ge0).
\]

\begin{proposition}[Value as a Schur inner product]\label{prop:reform}
With $a,b$ as above:
\begin{enumerate}
\item[\rm(A)] If $n\equiv0,1\pmod3$ then $n_1=n_2=b$ and
  \[
    G_\la(\zt)=\inner{s_\la,\;h_1^{a}\,\Theta^{b}}\ \in\ \ZZ_{\ge0}.
  \]
\item[\rm(B)] If $n\equiv2\pmod3$ then $n_1=n_2+1$, and writing
  $X=\inner{s_\la,h_1^{a}\Theta^{b}h_2}$, $Y=\inner{s_\la,h_1^{a}\Theta^{b}e_2}$ (both in
  $\ZZ_{\ge0}$),
  \[
    G_\la(\zt)=X+\zt\,Y .
  \]
  Consequently $G_\la(\zt)=0$ if and only if $X=Y=0$, and a sufficient condition for
  $G_\la(\zt)\neq0$ is
  \[
    X+Y=\inner{s_\la,\;h_1^{a+2}\,\Theta^{b}}>0 .
  \]
\end{enumerate}
\end{proposition}

\begin{proof}
By Lemma~\ref{lem:delta}, in case (A) $\delta=0$ so \eqref{eq:grouped} equals
$(h_1^2)^{n_0}\Theta^{b}$; multiplying by the prefactor $h_1^{e}$ gives
$h_1^{a}\Theta^{b}$, and pairing with $s_\la$ gives the formula. The inner product is a
nonnegative integer because $h_1^{a}\Theta^{b}$ is Schur-positive (a product of the
Schur-positive functions $h_1$ and $\Theta$, Lemma~\ref{lem:theta}).

In case (B) $\delta=1$, so \eqref{eq:grouped} equals
$(h_1^2)^{n_0}\Theta^{b}\phi=(h_1^2)^{n_0}\Theta^{b}(h_2+\zt\,e_2)$; with the prefactor
$h_1^{e}$ this is $h_1^{a}\Theta^{b}(h_2+\zt e_2)$, and pairing with $s_\la$ yields
$G_\la(\zt)=X+\zt Y$ with $X,Y$ as stated, again nonnegative integers. Since
$\{1,\zt\}$ is linearly independent over $\QQ$ and $X,Y\in\ZZ$, we have $G_\la(\zt)=0
\iff X=Y=0$. Finally $X+Y=\inner{s_\la,h_1^{a}\Theta^{b}(h_2+e_2)}
=\inner{s_\la,h_1^{a+2}\Theta^{b}}$ because $h_2+e_2=h_1^2$; as $X,Y\ge0$, $X+Y>0$
forces $(X,Y)\neq(0,0)$, hence $G_\la(\zt)\neq0$.
\end{proof}

Both cases thus rest on the single quantity $\inner{s_\la,h_1^{A}\Theta^{b}}$ with
$A+4b=n$ (here $A=a$ in case (A) and $A=a+2$ in case (B)). We characterise its
positivity.

\begin{lemma}[Pieri support criterion]\label{lem:pieri}
For $A\ge0$,
\[
  \inner{s_\la,\;h_1^{A}\Theta^{b}}=\sum_{\substack{\mu\subseteq\la\\ |\mu|=4b}}
       \inner{s_\mu,\Theta^{b}}\; f^{\la/\mu},
\]
where $f^{\la/\mu}=\#\SYT(\la/\mu)$. Hence $\inner{s_\la,h_1^{A}\Theta^{b}}>0$ if and only
if $\la$ contains some $\mu\in\supp(\Theta^{b})$ with $|\mu|=4b$.
\end{lemma}

\begin{proof}
Since $h_1^{A}=p_1^{A}$ and $s_\mu\,p_1^{A}=\sum_\la f^{\la/\mu}s_\la$ (the number of
standard skew tableaux of shape $\la/\mu$), writing
$\Theta^{b}=\sum_\mu\inner{s_\mu,\Theta^{b}}s_\mu$ and pairing with $s_\la$ gives the
displayed sum. All $\inner{s_\mu,\Theta^{b}}\ge0$ (Schur positivity) and
$f^{\la/\mu}>0\iff\mu\subseteq\la$. So the sum is positive iff some term is, i.e. iff
some $\mu\in\supp(\Theta^{b})$ of size $4b$ is contained in $\la$.
\end{proof}

\section{Full support of $\Theta^{b}$ for $b\ge3$}

\begin{lemma}[Removable strips]\label{lem:strip}
A partition $\rho$ admits a removable horizontal strip of size $4$ (i.e. some
$\mu\subseteq\rho$ with $\rho/\mu$ a horizontal $4$-strip) if and only if $\rho_1\ge4$,
and a removable vertical $4$-strip if and only if $\rho_1'\ge4$. Consequently, if
$|\rho|>9$ then $\rho$ admits a removable horizontal or vertical $4$-strip.
\end{lemma}

\begin{proof}
Taking $\mu=(\rho_2,\rho_3,\dots)$ makes $\rho/\mu$ a horizontal strip of size
$\sum_i(\rho_i-\rho_{i+1})=\rho_1$, and any intermediate $\mu_i\in[\rho_{i+1},\rho_i]$
keeps $\mu$ a partition with $\rho/\mu$ horizontal; thus removable horizontal strips of
every size $0,1,\dots,\rho_1$ exist, and none larger. The vertical statement is the
conjugate. If $\rho_1\le3$ and $\rho_1'\le3$ then $\rho\subseteq(3,3,3)$, so
$|\rho|\le9$; contrapositively $|\rho|>9$ forces $\rho_1\ge4$ or $\rho_1'\ge4$.
\end{proof}

\begin{lemma}[Full support]\label{lem:full}
For every $b\ge3$, $\supp(\Theta^{b})$ is all of $\{\mu:\mu\vdash 4b\}$.
\end{lemma}

\begin{proof}
\emph{Base $b=3$.} Expanding $\Theta^3$ in the Schur basis (a finite Pieri/LR
computation; all $77$ partitions of $12$ occur, see \S\ref{sec:verif}) shows
$\supp(\Theta^3)$ is full. \emph{Induction.} Assume $\supp(\Theta^{b})$ is full, $b\ge3$,
and let $\rho\vdash4(b+1)$. Then $|\rho|=4b+4\ge16>9$, so by Lemma~\ref{lem:strip}
$\rho$ has a removable horizontal or vertical $4$-strip; let $\mu\subseteq\rho$,
$\mu\vdash4b$, be the corresponding partition, so that $\rho/\mu$ is a horizontal (resp.
vertical) $4$-strip and hence the Littlewood--Richardson coefficient
$c^{\rho}_{\mu,(4)}=1$ (resp. $c^{\rho}_{\mu,(1^4)}=1$) by Pieri. Since $(4)$ and
$(1^4)$ lie in $\supp(\Theta)$ with positive coefficient and $\mu\in\supp(\Theta^{b})$
by hypothesis,
\[
  \inner{s_\rho,\Theta^{b+1}}=\sum_{\mu',\nu}\inner{s_{\mu'},\Theta^{b}}\inner{s_\nu,\Theta}
        c^{\rho}_{\mu'\nu}\ \ge\ \inner{s_{\mu},\Theta^{b}}\,\inner{s_{(4)\text{ or }(1^4)},\Theta}\,c^{\rho}_{\mu,\cdot}\ >0,
\]
all summands being nonnegative. Thus $\rho\in\supp(\Theta^{b+1})$, and
$\supp(\Theta^{b+1})$ is full.
\end{proof}

\begin{remark}
The hypothesis $b\ge3$ is sharp: $\supp(\Theta^{2})$ omits exactly $(6,1,1)$,
$(4,2,1,1)$, $(3,3,2)$, $(3,1^6)$, and $\supp(\Theta)=\{(4),(2,2),(1^4)\}$ omits
$(3,1),(2,1,1)$. This is the source of the two exceptional shapes below.
\end{remark}

\section{Main theorem}

\begin{theorem}\label{thm:main}
For every partition $\la$,
\[
  G_\la(\zt)=0 \iff \la\in\{(3,1),(2,1,1)\}.
\]
At these two shapes $G_\la(q)=q^5+q+1=\Phi_3(q)\,(q^3-q^2+1)$ and $\ord_{\zt}G_\la=1$.
For every other $\la$, $\ord_{\zt}G_\la=0$. In particular, if $\la$ has positive
$3$-weight (i.e. $\la$ is not a $3$-core) then $\ord_{\zt}G_\la=0$.
\end{theorem}

\begin{proof}
The set $\{n:b=n_2\le2\}=\{1,\dots,15\}\cup\{17\}$ is finite (it has maximum $17$; note
$n=16$ already has $n_2=3$); for every other $n$ one has $b\ge3$.

\emph{Case $b\ge3$ (i.e. $n\notin\{1,\dots,15,17\}$).} Let $\la\vdash n$ be arbitrary.
By Proposition~\ref{prop:reform} the nonvanishing of $G_\la(\zt)$ follows (in case (A)
directly, in case (B) via $X+Y$) from $\inner{s_\la,h_1^{A}\Theta^{b}}>0$ with
$A=n-4b\ge0$. By Lemma~\ref{lem:full}, $\supp(\Theta^{b})$ is full, so \emph{every}
$\mu\vdash4b$ lies in it; and $\la$, having $n=A+4b\ge4b$ boxes, contains some
$\mu\subseteq\la$ with $|\mu|=4b$ (remove $A$ removable corner boxes one at a time). By
Lemma~\ref{lem:pieri}, $\inner{s_\la,h_1^{A}\Theta^{b}}>0$. Hence $G_\la(\zt)\neq0$.

\emph{Case $b\le2$ (i.e. $n\in\{1,\dots,15,17\}$).} This is a finite computation
(\S\ref{sec:verif}): over all $\la\vdash n$ with $n\le17$, the only shapes with
$G_\la(\zt)=0$ are $(3,1)$ and $(2,1,1)$, both $3$-cores. Concretely, for the four shapes
of $n=4$, $\inner{s_\la,\Theta}$ vanishes only for $\la\in\{(3,1),(2,1,1)\}$
(Lemma~\ref{lem:theta}), and a check of cases (A)/(B) for the remaining
$n\in\{1,\dots,15,17\}$ exhibits a containing $\mu\in\supp(\Theta^{b})$ for every
$\la$.

Combining, $G_\la(\zt)=0$ iff $\la\in\{(3,1),(2,1,1)\}$. (At $n=4$, in case (A),
$\inner{s_\la,\Theta}$ is the coefficient of $s_\la$ in
$\Theta=s_{(4)}+2s_{(2,2)}+s_{(1^4)}$, which is zero precisely for the two shapes
$(3,1),(2,1,1)$ omitted from $\supp(\Theta)$.) The orders are read off from the
explicit polynomials. At $\la=(3,1)$ the three tableaux have descent sets
$\{1\},\{2\},\{3\}$ and weights $(w_1,w_2,w_3)=(1,0,5)$, so
\[
  G_{(3,1)}(q)=q^5+q+1=\Phi_3(q)\,(q^3-q^2+1),\qquad
  G_{(2,1,1)}(q)=q^6+q^5+q=q\,\Phi_3(q)\,(q^3-q+1),
\]
and in each case $\Phi_3$ does not divide the displayed cofactor (its value at $\zt$ is
$2-\zt^2\neq0$, resp. $2-\zt\neq0$), so $\ord_{\zt}G_{(3,1)}=\ord_{\zt}G_{(2,1,1)}=1$.
For all other $\la$, $G_\la(\zt)\neq0$ gives $\ord_{\zt}G_\la=0$. Since both exceptional
shapes are $3$-cores, every partition of positive $3$-weight has $\ord_{\zt}G_\la=0$.
\end{proof}

\section{Computational verification}\label{sec:verif}

All steps were checked by exact computation; scripts in
\texttt{\~{}/projects/scratch/} (\texttt{prove-master-id-check.py},
\texttt{prove-reformulation-check.py}, \texttt{prove-schur-support.py}).
\begin{itemize}
\item \textbf{Master identity (Prop.~\ref{prop:master}).} The character form
$G_\la(\zt)=\sum_{B\subseteq A}\chi^\la(2^{m-|B|}1^{2|B|+e})\prod_{k\in B}\frac{1+z_k}{2}
\prod_{k\notin B}\frac{1-z_k}{2}$ agrees with direct $\SYT$ evaluation of $G_\la(\zt)$
for all $\la\vdash n$, $n\le8$ ($66$ shapes, $0$ mismatches).
\item \textbf{Reformulation (Prop.~\ref{prop:reform}).} For all $\la\vdash n$, $n\le11$
($194$ shapes), the case-(A)/(B) values equal $G_\la(\zt)$ from direct enumeration, and
each of the quantities $G$ (case A), $X,Y$ (case B) is a nonnegative integer; $0$
exceptions.
\item \textbf{Lemma~\ref{lem:delta}.} $\delta\in\{0,1\}$ with $\delta=1\iff n\equiv2
\pmod3$ for all $n\le200$; and $\{n:n_2\le2\}=\{1,\dots,15,17\}$.
\item \textbf{Lemmas~\ref{lem:strip}, \ref{lem:full}.} Maximal removable horizontal
strip of $\rho$ equals $\rho_1$ (all $\rho$, $n\le10$); $\supp(\Theta^{3})$ and
$\supp(\Theta^{4})$ are full; $\supp(\Theta^{2})$ omits exactly the four listed shapes.
\item \textbf{Finite check (Theorem~\ref{thm:main}).} Over all $\la\vdash n$, $n\le17$
(case A via $\supp(h_1^{a}\Theta^{b})$ for $n\le16$; case B via $(X,Y)$ for
$n\in\{2,5,8,11,14,17\}$), the only $\la$ with $G_\la(\zt)=0$ are $(3,1),(2,1,1)$. This
is consistent with the independent exhaustive enumeration over $|\la|\le12$ recorded in
the order-law program.
\end{itemize}

\section{Remark: the obstruction for general $d\ge4$}

The master identity (Prop.~\ref{prop:master}) holds at any $\zeta_d$, and the same
grouping pairs conjugate factors into
\[
  (h_2+e_2\,\omega^{j})(h_2+e_2\,\omega^{-j})=h_2^2+c\,h_2e_2+e_2^2
   =s_{(4)}+(1+c)\,s_{(3,1)}+2\,s_{(2,2)}+(1+c)\,s_{(2,1,1)}+s_{(1^4)},
\]
with $\omega=\zeta_d$ and $c=2\cos(2\pi j/d)$. This quadratic factor is Schur-positive
precisely when $c\ge-1$, i.e. $|j|/d\le 1/3$. For $d=3$ the only paired frequency is
$j=1$, giving $c=-1$ exactly; the off-diagonal coefficient $1+c$ vanishes and one lands
on the clean $\Theta=s_{(4)}+2s_{(2,2)}+s_{(1^4)}$ that makes the whole argument run.
For $d\ge4$ some paired frequencies have $c<-1$, so the individual factors need not be
Schur-positive, and the reduction to a single nonnegative inner product breaks down. The
collapse $\ord_{\zeta_d}G_\la\in\{0,1\}$ is verified computationally for all $d\le8$ and
$|\la|\le12$, but a proof for general $d$ requires controlling the (possibly signed)
contributions of these factors; we leave this open.

\end{document}
